% =====================================================================
%  Betriebsanweisung (Gefahrstoff) EMAG EM-080 Universalreiniger
%  Eigenes PDF: Betriebsanweisung_Ultraschallreiniger.tex,
%  Seite in der Einweisung über \BAseite
%  Grundlage: Sicherheitsdatenblatt EMAG EM-080, überarbeitet am 15.05.2023:
%  Skin Irrit. 2 (H315), Eye Dam. 1 (H318), GHS05, Signalwort Gefahr,
%  P280, P305+P351+P338, pH 13,2 (Konzentrat) bzw. 11,1 (1 %), WGK 1.
%  Wird ein anderer Reiniger verwendet, diese BA an dessen SDB anpassen.
%  Das Sicherheitsdatenblatt muss im FabLab vorliegen.
% =====================================================================
\BAsetup{
  typ=gefahrstoff,
  titel={EMAG EM-080 Universalreiniger -- alkalisches Reinigungskonzentrat für Ultraschallbäder},
  kurztitel={Gefahrstoff: Ultraschallreiniger EMAG EM-080},
  taetigkeit={Dosieren des Konzentrats, Reinigen im Ultraschallbad, Entleeren, Lagern},
  nummer=BA-GS-02,
  stand=01.10.2026,
  verantwortlich={Name der verantwortlichen Person},
  signalwort=GEFAHR,
}

\BAkopf

% ---------- 1 Gefahren -----------------------------------------------
\BAblock{GEFAHREN FÜR MENSCH UND UMWELT}{GHS05}{%
  \begin{itemize}
    \item \textbf{Verursacht schwere Augenschäden} (H318). Schon ein Spritzer
          Konzentrat kann das Auge dauerhaft schädigen.
    \item Verursacht Hautreizungen (H315).
    \item Das Konzentrat ist eine starke Lauge (pH ca.\ 13), auch die
          gebrauchsfertige Lösung ist alkalisch (pH ca.\ 11).
    \item Reagiert heftig und unter Erwärmung mit konzentrierten Säuren.
    \item Schwach wassergefährdend (WGK 1).
  \end{itemize}}

% ---------- 2 Schutzmaßnahmen ----------------------------------------
\BAblock{SCHUTZMASSNAHMEN UND VERHALTENSREGELN}{M004,M009}{%
  \begin{itemize}
    \item Beim Umgang mit dem Konzentrat \textbf{Schutzbrille} und
          \textbf{Nitrilhandschuhe} tragen.
    \item Nur so viel dosieren wie nötig: 1--3\,\%, also 30--90\,ml auf
          3\,l Wasser. Mit Messbecher dosieren, nicht aus der Flasche gießen.
    \item Nicht mit anderen Reinigern, Säuren oder Lösemitteln mischen.
    \item Nicht in das Bad greifen, Teile nur im Korb herausnehmen und danach
          mit Wasser abspülen.
    \item Nach der Arbeit Hände waschen. Essen und Trinken verboten.
    \item Nur im Originalbehälter lagern, Flasche verschlossen halten,
          nie in Lebensmittelgefäße umfüllen.
  \end{itemize}}

% ---------- 3 Gefahrfall ---------------------------------------------
\BAblock{VERHALTEN IM GEFAHRFALL}{W001}{%
  \begin{itemize}
    \item Verschüttetes Konzentrat mit Papiertüchern oder Bindemittel
          aufnehmen, Fläche mit viel Wasser nachwischen, Betreuer informieren.
    \item Verschmutzte Kleidung sofort ausziehen.
    \item Das Produkt ist nicht brennbar. Bei einem Brand in der Umgebung
          mit Wasser, Schaum oder Sprühwasser löschen.
          \textbf{Feuerwehr: 112}
  \end{itemize}}

% ---------- 4 Erste Hilfe --------------------------------------------
\BAblock{ERSTE HILFE}{E003,E011}{%
  \begin{itemize}
    \item \textbf{Augen:} sofort 10--15 Min.\ bei geöffnetem Lid unter
          fließendem Wasser spülen, Kontaktlinsen entfernen, \textbf{Augenarzt}.
    \item \textbf{Haut:} sofort mit viel Wasser und Seife abwaschen.
    \item \textbf{Einatmen von Sprühnebel:} an die frische Luft, Arzt.
    \item \textbf{Verschlucken:} Mund ausspülen, reichlich Wasser trinken,
          kein Erbrechen herbeiführen, Arzt.
    \item Dem Arzt das Sicherheitsdatenblatt oder die Flasche zeigen.
  \end{itemize}\vspace{1.5mm}
  \textbf{Notruf: 112}\qquad Giftnotruf München: 089\,/\,19240\qquad
  FabLab: 09131\,/\,85\,28013}

% ---------- 5 Entsorgung (mit Fußzeile) --------------------------------
\BAletzterblock{SACHGERECHTE ENTSORGUNG}{}{%
  \begin{itemize}
    \item Gebrauchte, verdünnte Reinigungslösung mit reichlich Wasser in
          den Ausguss.
    \item \textbf{Konzentrat nicht in den Ausguss.} Reste und alte Flaschen
          an einen Betreuer geben, Entsorgung als Sonderabfall
          (Abfallschlüssel 20\,01\,29).
    \item Vollständig entleerte, ausgespülte Flaschen in die Wertstoffsammlung.
  \end{itemize}}
